Táto bakalárska práca sa zaoberá automatizáciou detekcie komponentov v historických rukou písaných nomenklátoroch pomocou metód umelej inteligencie. Nomenklátory sú šifrovacie systémy kombinujúce substitúciu písmen, n-gramov, klamáčov a kódov. Manuálna analýza týchto dokumentov je časovo náročná a nákladná, čo motivovalo výskum automatizovaných riešení. 

Cieľom práce bolo navrhnúť systém na detekciu častí šifrovacích kľúčov pomocou modelov YOLO, konkrétne verzie YOLO11. Práca zahŕňala prípravu dátovej sady pozostávajúcej z 89 digitalizovaných strán historických dokumentov, ich anotáciu do troch tried reprezentujúcich rôzne časti nomenklátorov, ako aj experimenty s rôznymi variantmi predspracovania dát. 

Výsledkom práce sú natrénované modely YOLO11 a ich vzájomné porovnanie. Analýza výsledkov prispieva k výskumu problematiky spracovania nomenklátorov. Detekcia jednotlivých častí nomenklátorov umožňuje ich následné automatizované spracovanie, čo prispieva k efektívnejšiemu výskumu histórie kryptografie. Zároveň poskytuje základ pre ďalšie vylepšenie metód spracovania historických rukopisov pomocou hlbokého učenia. 